\begin{frame}{Эволюция нижней границы информации Фишера (AMR-20)\footnote{G.A. Bochkin et al., \textit{Appl. Magn. Res.} \textbf{51}, 667-678, (2020)}}
\begin{columns}

    \column{0.6\textwidth}
    \begin{figure}
    \includegraphics[width=0.95\textwidth]{m2-by-time-in-zigzag-chain-with-n6-beta10.eps}
    %\caption{Hанопора со спин-несущих молекулами во внешнем сильном магнитном поле $\vec B$}
    \end{figure}

    \column{0.4\textwidth}
    Зависимость нижней границы квантовой информации Фишера
    $$ F_Q = 2M_2(\tau, T) $$
    от безразмерного времени $D_1 \tau$
    для шести спинов
    при температуре $T = 2.5 \times 10^{-3}$ $(b = 10)$.
    Область многочастичной запутанности ограничена горизонтальными линиями $k = 1$, $k = 5$.
\end{columns}
\end{frame}
\note{
   Мы не можем точно сказать сколько спинов запутанно,
   но можем дать нижнюю оценку на количество запутанных между собой спинов.
}


\begin{frame}{Максимальное количество запутанных спинов (AMR-20)}
\begin{columns}

    \column{0.5\textwidth}
    \begin{figure}
    \includegraphics[width=\textwidth]{nent-by-n-b10.eps}
    %\caption{Hанопора со спин-несущих молекулами во внешнем сильном магнитном поле $\vec B$}
    \end{figure}

    \column{0.5\textwidth}
    Зависимость максимального количества запутанных спинов $N_\mathrm{ent}$ от длинны цепи при температуре $T = 2.5 \times 10^{-3}$ $(b = 10)$.

    \vspace{0.5cm}

    \alert{Создание запутанных кластеров в рассматриваемых зигзагообразных цепочках ограничено слабыми дипольными взаимодействиями удаленных спинов}.
\end{columns}
\end{frame}
\note{
    Установлено что, при фиксированной температуре с ростом числа спинов в цепи размер запутанных кластеров может уменьшаться;
}


\begin{frame}{Максимальное количество запутанных спинов (AMR-20)}
  \begin{columns}

     \column{0.6\textwidth}
     \begin{figure}
     \includegraphics[width=\textwidth]{nent-by-beta-n8-n6.eps}
     \end{figure}

     \column{0.4\textwidth}
     Зависимость максимального количества запутанных спинов $N_\mathrm{ent}$ от обратной температуры $b$
     $$ b=\dfrac{\hbar\omega_0}{kT} $$
     в зигзагообразной цепочки из 6 и 8 спинов.

  \end{columns}
\end{frame}
\note{
    Для 6 спинов накачка происходит быстрее, чем для 8 спинов.
}
